\section{Implementation and verification}
\label{ch:implementation}

Part II identifies the traffic, equilibrium, and policy objects required by
the welfare derivative. An application must represent those objects in common
units, construct a coherent baseline, and preserve the policy definition
through the numerical calculation. Data and calibration come first because
the validity of later software checks depends on that representation.

\subsection{Data and calibration}
\label{ch:data}

Raw transport records rarely coincide with the model's nodes, flows, or cost
units. This subsection states the model-ready schemas, accounting identities,
parameter sources, and policy fields that convert those records into inputs
the derivative can use.

\subsubsection{Model-ready inputs}

The theory identifies the inputs needed by the derivative; the data contract
states how an application must represent them.
The generic adapter reads two CSV files and one TOML configuration. The files
are deliberately narrow. They contain a baseline that already satisfies the
selected model's accounting identities; they are not a warehouse for raw
transport records.

For economic geography, the node file contains:
\begin{lstlisting}
node_id,labor,income,longitude,latitude
\end{lstlisting}
Labor and income shares must be finite, strictly positive, and normalized.
Zero-mass locations are outside the interior baseline accepted by the generic
loader. Coordinates are optional for analysis but required for maps.

For urban commuting, the node file instead contains:
\begin{lstlisting}
node_id,residents,employment,longitude,latitude
\end{lstlisting}
Residence and workplace margins must be finite, positive, and normalized.
Their geographic distributions may differ even though their aggregate totals
refer to the same commuter population.

The edge-mode file contains one row for each active mode on each directed edge:
\begin{lstlisting}
edge_id,physical_link_id,origin,destination,mode,flow,
origin_terminal_id,destination_terminal_id,
congestion_elasticity
\end{lstlisting}
The terminal and edge-specific congestion columns are optional unless selected
by the configuration. The pair \((\code{edge\_id},\code{mode})\) must be
unique. One directed \code{edge\_id} may therefore appear once for each active
mode, but all such rows must share endpoints and a physical-link identifier.
Each endpoint must exist in the node table.

\implementationbox{\textbf{Recovering baseline cost derivatives from flows.}
Given the recursive-route and modal-logsum assumptions, consistent
edge--mode flows and node margins reveal the continuation and modal shares
needed for the local cost derivative. The package reconstructs those shares
without separately observed route-level generalized costs. This identification
also requires calibrated route and modal elasticities and flows that satisfy
the model's accounting identities. Subsection
\ref{sec:object-status} distinguishes observed, reconstructed, calibrated, and
computed quantities.}

\subsubsection{Flow concepts and units}

Economic geography uses value traffic normalized by world income. The urban
model uses commuter traffic normalized by total commuters. Vehicle counts,
passenger trips, tons, and scheduled services are not automatically
interchangeable with either flow concept. A preprocessing pipeline may convert
a physical measure to the model's flow, but that conversion must be explicit
and defensible.

All active modal flows on an edge sum to aggregate edge traffic:
\[
 \Xi_{ij}=\sum_m\Xi_{ij,m},
 \qquad
 s_{ij,m}=\frac{\Xi_{ij,m}}{\Xi_{ij}}.
\]
The package constructs modal shares from the input baseline. It does not use
\(\eta\) to overwrite those baseline shares. The elasticity controls the local
response after a shock.

For economic geography, labor and income shares each sum to one, and edge value
flows use the same world-income denominator. For urban commuting, residence and
workplace margins each sum to one, and edge flows use the same total-commuter
denominator. In either model, modal shares sum to one within a directed edge.
Edge traffic does not generally sum to one because the same shipment or
commuter can traverse several edges along a route.
Renormalizing the edge table to unit total would therefore change the welfare
scale and violate the route accounting.

The unit contract extends to congestion. A log elasticity such as
\(\gamma=0.092\) is invariant to proportional rescaling of a correctly
matched flow measure. It is not invariant to replacing value or commuter
traffic with an unrelated vehicle, passenger, or schedule measure.

\subsubsection{Accounting checks}

For economic geography, outgoing and incoming traffic must recover the same
trade exposure:
\[
 \mathcal T_i^G=\frac{O_i}{1-s_i^x}
               =\frac{I_i}{1-s_i^y}.
\]
For urban commuting, residence and workplace absorption must recover the same
commuting exposure:
\[
 \mathcal T_i^U=l_i^F+O_i=l_i^R+I_i.
\]
Using the identity selected by the configuration, the loader first verifies
node and mode identifiers, positivity, and finiteness. It then checks flow
conservation and agreement between alternative measures of the exposure stock.
For the requested policy, it also verifies reciprocal physical-link pairing,
coverage of the policy mode, terminal identifiers for modes subject to terminal
congestion, and congestion keys against the observed active modes.

Perform any symmetrization, padding, modal rescaling, or balancing before
loading the generic adapter. When empirical sources do not agree, create a
preprocessing artifact that states the transformation and supply its output to
the package. The run manifest records the hashes of the transformed inputs.

\subsubsection{Preparing model-ready flows}

Different sources require different preprocessing. This table identifies the
main empirical gap without presuming a universal conversion.
\begin{longtable}{L{0.34\textwidth}L{0.57\textwidth}}
\toprule
Available information & Required construction \\
\midrule
\endhead
Directed edge value flows &
Apply unit, price-basis, and accounting checks directly. \\
Vehicle or passenger counts &
Map the physical count to the model's value or commuter flow and match the
congestion elasticity to that measure. \\
Origin--destination value flows &
Assign each OD flow to directed edge--modes under a declared routing model. \\
Scheduled transit service &
Combine schedules with ridership or a modeled commuter baseline; schedules
alone are not observed traffic. \\
Imports and exports with different margins &
Project them onto declared common margins while preserving and reporting the
chosen support. \\
Incomplete modal coverage &
Define a defensible residual mode or limit the application to the observed
choice set. \\
Zero-flow proposed infrastructure &
Use a finite counterfactual or entry model; the interior local derivative does
not identify the new alternative. \\
\bottomrule
\end{longtable}

Reconcile inconsistent source matrices in a separate, versioned balancing
step. First choose the target node margins and preserve
structural zeros that represent unavailable connections. Then solve a
constrained projection, such as iterative proportional fitting when its
support conditions hold, that changes the source matrix as little as possible
while matching those margins. Report the projection rule, fixed margins,
support changes, convergence tolerance, iteration count, maximum margin error,
and a distance between raw and balanced flows. Archive both the raw-source
hashes and the model-ready output hash. The freight adapter below implements
this procedure for directional port trade; other applications must justify
their own targets and support.

\subsubsection{Parameters}

After the accounting inputs are fixed, the configuration supplies the
elasticities governing spatial, route, modal, and congestion responses. The
economic-geography parameter vector is
\[
 (\alpha,\beta,\sigma,\eta,\gamma_{e,m}).
\]
Here \(\gamma_{e,m}\) denotes the declared edge--mode congestion elasticity;
the route curvature is tied to \(\sigma-1\) by the paper's theorem. The paper
baseline uses
\[
 \alpha=0.10,\quad
 \beta=-0.30,\quad
 \sigma=9,\quad
 \eta=1.099,\quad
 \gamma_{\mathrm{road}}=0.092.
\]
These values are external calibration inputs, not parameters estimated by the
package. The spatial elasticities \(\alpha=0.10\) and \(\beta=-0.30\) follow
the economic-geography calibration in \citet{AllenArkolakis2014}, as used in
the December 2020 revision of \citet{AllenArkolakis2020}. The route curvature
\(\sigma-1=8\) matches the value used by
\citet{FuchsWong2026Multimodal}, while their estimated modal elasticity gives
\(\eta=1.099\). The road-congestion elasticity
\(\gamma_{\mathrm{road}}=0.092\) comes from
\citet{AllenArkolakis2020}. The endpoint-terminal extension uses the
\(0.096\) dwell-time elasticity estimated by
\citet{FuchsWong2026Multimodal}. These citations identify the source of each
baseline value; they do not substitute for sensitivity analysis in a new
application.

The urban parameter vector replaces \(\sigma\) with commuting heterogeneity
\(\theta_U\):
\[
 (\alpha,\beta,\theta_U,\eta,\gamma_{e,m}).
\]
The package uses \(\theta_U\) as the urban route curvature. Its modal and
congestion elasticities remain application inputs. The Seattle candidate
transfers \(\eta=1.099\) from the freight application and reports sensitivity
rather than presenting it as a Seattle estimate.

Sensitivity should vary one parameter while remaining on a regular branch.
The paper sensitivity runner maintains
\[
 \alpha+\beta<0,\qquad e<0,\qquad
 |e|\geq0.05,\qquad 1+\alpha+\beta>0.05.
\]
Thus ``net dispersion'' means \(\alpha+\beta<0\), and the remaining conditions
keep the path on the baseline sign branch and away from the two transformation
singularities. The generic model builder requires the corresponding
denominators to be nonzero. Both models reject singular route, transport, or
spatial systems, and the economic-geography implementation rejects an
independent route curvature not covered by the theorem.

\subsubsection{Policy definition}

The configuration states the policy mode and whether the output concerns
directed arcs, reciprocal physical links, or both. It also records the
proportional shock used to translate elasticities into reported gains and
indicates whether the policy combines several edge--mode shocks in a bundle.

A policy bundle is a sparse list:
\begin{lstlisting}
bundle_id,edge_id,mode,weight
\end{lstlisting}
The first-order welfare effect is linear in these weights. Duplicate entries,
unknown edges, inactive modes, and nonfinite weights are rejected.

\subsubsection{External-data adapters}

The freight replication adapter, retained in code under its historical
\code{RSUE} name, performs transformations required to reproduce the paper's
research pipeline. Unlike the generic adapter, it may pad, symmetrize, filter,
or rescale because those operations are part of the declared replication
contract. Each transformation is listed in TOML and written to the manifest.

The Seattle adapter combines a published urban archive with LODES commuting,
ACS mode shares, and a pinned historical GTFS network
\citep{CensusLODES,CensusACS2017B08301,KingCountyGTFS2017}. The raw sources
remain outside Git. Schedule data identify service paths and frequency, not
observed route ridership. The model builder therefore reports its construction
and external comparisons without presenting schedule counts as passenger
flows.

\warningbox{A successful schema check establishes internal consistency, not
external validity. Validation against untargeted traffic, ridership, prices,
or location outcomes is a separate empirical exercise.}

Once the data and policy satisfy this contract, the computation constructs the
route basis, assembles the relevant equilibrium system, and converts each
policy forcing into a welfare effect. The next subsection follows that
sequence.
